% Shared thesis preamble.  Paper-specific commands retained below allow the
% source manuscripts to be incorporated with minimal textual intervention.
\usepackage[pdftex,dvipsnames]{xcolor}
\usepackage[
  a4paper,
  inner=35mm,
  outer=25mm,
  top=25mm,
  bottom=30mm,
  headheight=15pt
]{geometry}
\usepackage{setspace}
\onehalfspacing
\usepackage{fancyhdr}
\pagestyle{fancy}
\fancyhf{}
\fancyhead[R]{\thepage}
\fancyhead[L]{\nouppercase{\leftmark}}
\renewcommand{\headrulewidth}{0.3pt}
\usepackage{emptypage}


% Recommended, but optional, packages for figures and better typesetting:
\usepackage{microtype} %  Use "microtype" always; just \usepackage{microtype} and it'll make stuff awesome for free
\usepackage{graphicx}
\usepackage{subcaption}
% \usepackage{subfigure}
\usepackage{verbatim}
\usepackage{booktabs} % READ the documentation of "booktabs" and then use it whenever you are preparing tables (the manual is awesome) https://ctan.org/pkg/booktabs
\usepackage{tabularx}
\usepackage{longtable}
\usepackage{array}
\usepackage{algpseudocode}
\usepackage{amsmath}
\usepackage{algorithm}
\usepackage{mathpazo}
\usepackage{caption}
\usepackage{wrapfig}
\usepackage{mdframed}
\usepackage{mathtools}
\usepackage{float}
\newcommand{\squeezeup}{\vspace{-30mm}}
\usepackage{pythonhighlight}




\usepackage{natbib}

% hyperref makes hyperlinks in the resulting PDF.
% If your build breaks (sometimes temporarily if a hyperlink spans a page)
% please comment out the following usepackage line and replace
% \usepackage{icml2024} with \usepackage[nohyperref]{icml2024} above.
\usepackage[hyphens]{url}
\usepackage[
colorlinks=true,
linkcolor=blue,
urlcolor=blue,
citecolor=blue,
anchorcolor=blue,
pdftitle={Economic Theory for AI Alignment},
pdfauthor={Abhimanyu Pallavi Sudhir}
]{hyperref}
\usepackage{siunitx}
\usepackage{enumitem}
\usepackage{xfrac}
\usepackage{changepage}

%\usepackage[accepted]{icml2024}
%\usepackage{icml2024}

\usepackage{ifthen}
\usepackage{amssymb}% http://ctan.org/pkg/amssymb
\usepackage{pifont}% http://ctan.org/pkg/pifont
\newcommand{\xmark}{\ding{55}}%

% Attempt to make hyperref and algorithmic work together better:
\newcommand{\theHalgorithm}{\arabic{algorithm}}
\makeatletter
\providecommand{\theHALG@line}{}
\renewcommand{\theHALG@line}{\thealgorithm.\arabic{ALG@line}}
\makeatother

\usepackage{threeparttable}



\usepackage{xargs}

\usepackage{xparse}
\usepackage{amsfonts}
\usepackage{amsmath}
\usepackage{amssymb}
\usepackage{amsthm}

% DO NOT USE \usepackage{geometry} IN THE ICML TEMPLATE

\usepackage{pgfplots}
\pgfplotsset{compat=1.17}
\usetikzlibrary{calc}

\newlength\tindent
\setlength{\tindent}{\parindent}
\setlength{\parindent}{1.5em}
\setlength{\parskip}{0.25em}
\renewcommand{\indent}{\hspace*{\tindent}}

% math
\NewDocumentCommand{\Indicator}{m}{\ensuremath{\delta_{#1}}}
\NewDocumentCommand{\Reals}{}{\mathbb{R}}
\NewDocumentCommand{\bools}{}{\{\top,\bot\}}
\NewDocumentCommand{\probs}{}{[0,1]}
\NewDocumentCommand{\none}{}{\mathrm{None}}
\NewDocumentCommand{\boolsopt}{}{\{\top,\bot,\none\}}
\newcommand{\eps}{\varepsilon}
\newcommand{\abs}[1]{\lvert #1 \rvert}
\newcommand{\norm}[1]{\lVert #1 \rVert}
\DeclareMathOperator*{\argmax}{arg\,max}
\DeclareMathOperator*{\argmin}{arg\,min}

% forecasting question sets
\NewDocumentCommand{\ForecastingQuestions}{}{\operatorname{Prop}}
\NewDocumentCommand{\Resolutions}{}{\Theta}
\NewDocumentCommand{\ResolutionsOpt}{}{\Theta'}
\NewDocumentCommand{\Forecasts}{}{\operatorname{Forecast}}
\NewDocumentCommand{\Resolution}{}{\theta}
\NewDocumentCommand{\ResolutionsVec}{}{\Omega}
\NewDocumentCommand{\resolutionvec}{}{\omega}

% forecasters
\NewDocumentCommand{\Forecaster}{}{\mathbb{F}}
\NewDocumentCommand{\Prob}{}{\mathbf{P}}
\NewDocumentCommand{\finset}{}{\operatorname{finset}}
\newcommand{\citeps}{\citep}
\newcommand{\citets}{\citet}

% tuples
\NewDocumentCommand{\Relation}{}{\mathcal{R}}
\NewDocumentCommand{\CCheck}{}{\mathcal{S}}
\NewDocumentCommand{\Violation}{}{\mathcal{V}}
\NewDocumentCommand{\Arbitrageur}{}{\mathcal{A}}
\NewDocumentCommand{\FullInstantiator}{}{\mathcal{J}}
\NewDocumentCommand{\Instantiator}{}{\mathcal{I}}
\NewDocumentCommand{\Generator}{}{\mathcal{G}}

% daniel's frequentist approach
\NewDocumentCommand{\TrueProb}{}{\mathbb{T}}
\NewDocumentCommand{\Normal}{}{\mathsf{N}}

% consistent forecaster
\newcommand{\CFCaster}{\texttt{ArbitrageForecaster}}
\NewDocumentCommand{\CFC}{o m}{%
  \langle #2 \rangle%
  \IfValueT{#1}{_{#1}}%
}

% Todo notes
\setlength {\marginparwidth }{2cm} % !!! REMEMBER TO COMMENT THIS OUT WHEN DONE WITH TODONOTES !!!
\usepackage[
    disable, % UNCOMMENT TO TURN OFF TODONOTES
    textsize=tiny,
]{todonotes}

% Inline comment commands
% unimp is if it's not crucial
\newcommandx{\dpaleka}[2][1=]{\todo[linecolor=cyan,backgroundcolor=cyan!25,bordercolor=cyan,#1]{DP: #2}}
\newcommandx{\dpalekaunimp}[2][1=]{\todo[linecolor=cyan!50,backgroundcolor=cyan!12.5 bordercolor=cyan!50,#1]{DP: #2}}
\newcommandx{\manyu}[2][1=]{\todo[linecolor=black,backgroundcolor=black!25,bordercolor=black,#1]{APS: #2}}
\newcommandx{\manyuunimp}[2][1=]{\todo[linecolor=black!50,backgroundcolor=black!12.5,bordercolor=black!50,textcolor=black!50,#1]{APS: #2}}
\newcommandx{\Alejandro}[2][1=]{\todo[linecolor=OliveGreen,backgroundcolor=OliveGreen!25,bordercolor=OliveGreen,#1]{Al: #2}}
\newcommandx{\Alejandrounimp}[2][1=]{\todo[linecolor=OliveGreen!50,backgroundcolor=OliveGreen!12.5,bordercolor=OliveGreen!50,#1]{Al: #2}}
\newcommandx{\towrite}[2][1=]{\todo[linecolor=black,backgroundcolor=black,bordercolor=black,textcolor=white#1,inline]{ADD: #2}}
\newcommandx{\vin}[2][1=]{\todo[linecolor=magenta,backgroundcolor=magenta!25,bordercolor=magenta,#1]{VB: #2}}
\newcommandx{\vinunimp}[2][1=]{\todo[linecolor=magenta!50,backgroundcolor=magenta!12.5,bordercolor=magenta!50,#1]{VB: #2}}
\newcommandx{\adam}[2][1=]{\todo[linecolor=blue,backgroundcolor=blue!25,bordercolor=blue,#1]{Adam: #2}}
\newcommandx{\adamunimp}[2][1=]{\todo[linecolor=blue!50,backgroundcolor=blue!12.5,bordercolor=blue!50,#1]{Adam: #2}}
\newcommandx{\evan}[2][1=]{\todo[linecolor=red,backgroundcolor=red!25,bordercolor=red,#1]{Evan: #2}}
\newcommandx{\evanunimp}[2][1=]{\todo[linecolor=red!50,backgroundcolor=red!12.5,bordercolor=red!50,#1]{Evan: #2}}
\newcommandx{\ft}[2][1=]{\todo[linecolor=orange,backgroundcolor=orange!25,bordercolor=orange,#1]{FT: #2}}



% \newcommandx{\thiswillnotshow}[2][1=]{\todo[disable,#1]{#2}}


\newcommand{\nicett}[1]{{\fontfamily{lmvtt}\fontsize{10.5pt}{12pt}\selectfont #1}}
%\newcommand{\gptturbo}{\nicett{gpt-3.5-turbo}}


% Use the following line for the initial blind version submitted for review:
%\usepackage{icml2024}

% If accepted, instead use the following line for the camera-ready submission:

% if you use cleveref..
\usepackage[capitalize,noabbrev]{cleveref}


%%%%%%%%%%%%%%%%%%%%%%%%%%%%%%%%
% THEOREMS
%%%%%%%%%%%%%%%%%%%%%%%%%%%%%%%%

\theoremstyle{plain}
\newtheorem{theorem}{Theorem}[section]
\newtheorem{proposition}[theorem]{Proposition}
\newtheorem{lemma}[theorem]{Lemma}
\newtheorem{corollary}[theorem]{Corollary}
\theoremstyle{definition}
\newtheorem{definition}[theorem]{Definition}
\newtheorem{example}[theorem]{Example}
\newtheorem{assumption}[theorem]{Assumption}
\newtheorem{nontheorem}[theorem]{Non-theorem}
\theoremstyle{remark}
\newtheorem{remark}[theorem]{Remark}
\newtheorem*{notation*}{Notation}

%%%%%%%%%%%%%%%%%%%%%%%%%%%%%%%%
% FORMATTING
%%%%%%%%%%%%%%%%%%%%%%%%%%%%%%%%

% this some nice formatting taken from https://arxiv.org/abs/2404.13208
\definecolor{pastelred}{HTML}{FFC0C0}
\definecolor{pastelorange}{HTML}{FFDAC0}
\definecolor{pastelyellow}{HTML}{FFFDC0}
\definecolor{pastelgreen}{HTML}{a0e7a0}
\definecolor{pastelcyan}{HTML}{C0FFFD}
\definecolor{pastellightblue}{HTML}{D0E0FF}
\definecolor{pastelblue}{HTML}{C0C0FF}
\definecolor{pastelviolet}{HTML}{DAC0FF}
\definecolor{pastelmagenta}{HTML}{FFC0FF}
\definecolor{pastelrose}{HTML}{FFC0DA}

\definecolor{darkred}{HTML}{C23B22}
\definecolor{darkgreen}{HTML}{1cc650}
\definecolor{lightgreen}{HTML}{caee9c}

\usepackage{tcolorbox}
\tcbuselibrary{breakable,skins}

\newenvironment{promptbox}{
  \begin{figure}[hbt!]
  \begin{tcolorbox}[left=1.5mm, right=1.5mm, top=1.5mm, bottom=1.5mm]
  \raggedright
  \small
  \textbf{User:} \texttt{...} \\[4pt]
  \textbf{Assistant:} \texttt{...} \\[4pt]
}{
  \end{tcolorbox}
  \end{figure}
}

\newenvironment{emptypromptbox}[1][] % [Title (optional)]
{
  \begin{tcolorbox}[left=1.5mm, right=1.5mm, top=1.5mm, bottom=1.5mm]
    \raggedright
    \small
    \ifx\relax#1\relax\else
      \begin{center}
        {\normalsize \textbf{\color{black} #1}}
      \end{center}
    \fi
}{
  \end{tcolorbox}
}


\newtcolorbox{questionbox}[2][]{colback=gray!10!white, colframe=gray!40!black, title=#2, #1}


%\newcommand{\NegChecker}{\textsc{NegChecker}}
%\newcommand{\ConsequenceChecker}{\textsc{ConsequenceChecker}}
%\newcommand{\ParaphraseChecker}{\textsc{ParaphraseChecker}}
%\newcommand{\AndChecker}{\textsc{AndChecker}}
%\newcommand{\ButChecker}{\textsc{ButChecker}}
%\newcommand{\CondChecker}{\textsc{CondChecker}}
%\newcommand{\OrChecker}{\textsc{OrChecker}}
%\newcommand{\AndOrChecker}{\textsc{AndOrChecker}}
%\newcommand{\CondCondChecker}{\textsc{CondCondChecker}}
%\newcommand{\SymmetryAndChecker}{\textsc{SymmetryAndChecker}}
%\newcommand{\SymmetryOrChecker}{\textsc{SymmetryOrChecker}}


\newcommand{\NegChecker}{\textsc{Negation}}
\newcommand{\ConsequenceChecker}{\textsc{Consequence}}
\newcommand{\ParaphraseChecker}{\textsc{Paraphrase}}
\newcommand{\AndChecker}{\textsc{And}}
\newcommand{\ButChecker}{\textsc{But}}
\newcommand{\CondChecker}{\textsc{Cond}}
\newcommand{\OrChecker}{\textsc{Or}}
\newcommand{\AndOrChecker}{\textsc{AndOr}}
\newcommand{\CondCondChecker}{\textsc{CondCond}}
\newcommand{\SymmetryAndChecker}{\textsc{SymmetryAnd}}
\newcommand{\SymmetryOrChecker}{\textsc{SymmetryOr}}
\newcommand{\ExpectedEvidenceChecker}{\textsc{ExpEvidence}}

\newcommand{\Yes}{\textsc{Yes}}
\newcommand{\No}{\textsc{No}}
\newcommand{\None}{\textsc{N/A}}

\newcommand{\gptfouro}{\texttt{gpt-4o}}
\newcommand{\gptfouromay}{\texttt{gpt-4o-2024-05-13}}
\newcommand{\gptfouroaug}{\texttt{gpt-4o-2024-08-06}}
\newcommand{\gptfouromini}{\texttt{gpt-4o-mini}}
\newcommand{\gptfourominijul}{\texttt{gpt-4o-mini-2024-07-18}}
\newcommand{\claudesonnet}{\texttt{claude-3.5-sonnet}}
\newcommand{\claudehaiku}{\texttt{claude-3-haiku}}
\newcommand{\gptthreefiveturbo}{\texttt{gpt-3.5-turbo}}
\newcommand{\llamaeight}{\texttt{llama-3.1-8B}}
\newcommand{\llamaseventy}{\texttt{llama-3.1-70B}}
\newcommand{\llamafourhundred}{\texttt{llama-3.1-405B}}
\newcommand{\perplexityhuge}{\texttt{perplexity/llama-3.1-sonar-huge-128k-online}}
\newcommand{\perplexitylarge}{\texttt{perplexity/llama-3.1-sonar-large-128k-online}}
\newcommand{\oonemini}{\texttt{o1-mini}}
\newcommand{\oonepreview}{\texttt{o1-preview}}
\newcommand{\textembedsmall}{\texttt{text-embedding-3-small}}
\newcommand{\verbatimbox}[1]{{\obeyspaces\obeylines\detokenize{#1}}}

% Thesis metadata. Update this in one place before submission.
\newcommand{\thesissubmissiondate}{August 2026}

% Every sentence newly composed to connect source manuscripts must be wrapped
% in this environment.  It is deliberately visible in both source and PDF.
\newenvironment{drafttext}{%
  \begin{tcolorbox}[
    enhanced,
    breakable,
    colback=yellow!8,
    colframe=Goldenrod!65!black,
    boxrule=0.5pt,
    left=2mm,right=2mm,top=1mm,bottom=1mm,
    title={\footnotesize\sffamily DRAFT GLUE --- AUTHOR REVIEW REQUIRED},
    fonttitle=\bfseries
  ]
}{%
  \end{tcolorbox}
}

\newcommand{\chapterprovenance}[1]{%
  \begin{quote}\small\itshape\textbf{Source note.} #1\end{quote}
}

% Helpers emitted by Pandoc for the Markdown theory manuscript.
\providecommand{\tightlist}{%
  \setlength{\itemsep}{0pt}\setlength{\parskip}{0pt}}
\providecommand{\real}[1]{#1}

\setcounter{secnumdepth}{3}
\setcounter{tocdepth}{2}

% The consistency manuscript uses one level below subsubsection.
\newcommand{\subsubsubsection}[1]{\paragraph{#1}\mbox{}\\}

% These appear in a figure caption and therefore must exist when the list of
% figures is read, before the appendix that normally defines them.
\NewDocumentCommand{\priceP}{}{\Forecaster(P)}
\NewDocumentCommand{\priceQ}{}{\Forecaster(Q)}
